\chapter*{In Solomon's Porch}
\textbf{"Give me new consolation, great and strong, of which I nave never
heard or read."}--\textit{Pliny}.

With increasing frequency Christ was now speaking his prophecies of
the life immortal.  In his earlier ministry he had been dwelling
upon the presence of the divine kingdom in the earth, the practical
conditions for membership therein, and the inclusion of Gentile as
well as Jew in the gracious provision.  Novel were his words.
Whoever had heard his discourse on the Mount or the parable of the
lost sheep was rich beyond the modern sons of men.  But now, in the
closing period of his stay with mortals, he was more frequently
foretelling the life to come.  Like a footworn traveler drawing
near the homeland, he was keenly anticipating his return to the
spirit world.  Those who listened to him heard majestic intimations
of a celestial country which eye had not beheld.  Nor is it to be
thought that the Gospels, in their restricted pages, have recorded
half his words concerning the heavenly land.

Now comes the opportunity for Quintus himself to hear this new
Teacher of the Jews.  A messenger from Pilate, sent on an errand to
the headquarters at Scopus, brings the tidings that Christ is in
Jerusalem as a visitor at the Feast of Dedication.  Favored are
those who hear through the years the world's commanding voices;
beyond estimate is the high privilege now granted Quintus.

"I will hasten in to Hierosolyma," he says to Aulus, who is
detained by camp duties; "I will hear him for myself; and I will
bring you back report as to this latest prophet of immortality."


With his soldier's cloak about him, in protection against the
winter's chill, Quintus is away to Jerusalem.  The national Feast
of Dedication attracts his notice.  A courteous Hebrew explains to
him that the joyful festival commemorates the cleansing of the
Temple after its profanation by Antiochus Epiphanes, two hundred
years before.  The procession of pious Jews, carrying their palm
branches and marching to the heights of Moriah, the chanting of the
great Hallel within the imposing fane, the ascription of praise to
Jehovah all impress the keen-eyed soldier.

The enthusiasm of it all!  Though of other blood, Quintus clearly
feels the thrill of patriotism that stirs the multitude about him;
and he understands in some measure their impatient waiting for the
coming prince who shall deliver Israel.


But is this all?  Instead it is only the beginning of the wonders
which the serious Quintus is to witness.  Forth he passes to the
eastern cloister of the Temple, known then among the Jews as
Solomon's Porch, in memory of their illustrious king.  The
bystanders tell Quintus that it is built of a fragment of the first
Temple which Nebuchadnezzar had left standing.  As the soldier
looks down the far-reaching aisle, he sees a quadruple row of white
Corinthian columns, one hundred and sixty in number, and extending
a length of many hundred feet.  The vista is most amazing.
Accustomed though he has been all his days to the magnificence of
the Roman architecture, he yields in willing admiration to the
splendors of the Solomonic porch.

Then--he sees the Christ!  Walking through that forest of massive
columns is the superlative Jew of his times, and of all times.  For
now--when the voices of that winter day are still, and Solomon's
Porch has vanished where stood those blessed feet--there is no
earthly measurement by which to estimate the Man whom Quintus saw.


Among the throng that surround him hostile Pharisees challenge him
to tell them plainly if he be the foretold Messiah.  With impatient
hearts they have waited long for their redemption.  Let him say if
their deliverer has now come.  Then shall they throw off the yoke
of the detested Roman rule and renew their ancient monarchy with
enlarging influence and increasing splendors.

Memorable words in answer does Quintus hear.  The Stranger  puts
aside  the thought of the Jewish struggle for an earthly throne,
and turns in his fancy to the quiet pastures where feed the flocks.
He is a guardian Shepherd; Israel and all the world besides are his
cherished sheep.  Those who are truly his shall hear his guiding
voice, and shall follow him.  They shall never perish.  From the
hand of the Shepherd no vandal shall steal his own away.  How the
words thrill!  Sometimes Quintus has seen in the Judaean pastures
the keeper with his flocks, and knows how unchanging is his
fidelity.  It is as if this watcher in his devotion is anticipating
the faithfulness of the greater Shepherd.  How entrancing is the
lesson to this seeking soldier from beyond the Adriatic!

Then does the Christ add another word more surprising than the
rest.  To men who are his sheep he makes a promise that compasses
the furthest limit of the eternities.  Of such he says: "Unto those
who follow me I will give the Life of the Ages.  Beyond the tomb
they are to live on forevermore."  Nor to the Jews alone, amid the
maze of those Corinthian columns, does the coming Shepherd speak.
The listening Roman soldier, wearing the armor of the empire on the
Tiber, comes within the circle of his promise.  Into the face of
Quintus he looks and benignly says: "There are other sheep not of
the Jewish pasture, to whom I shall give this unending life.  I
covet your great empire as my own.  O soldier of the Caesars,
follow after me!"


Back to the camp on Scopus the soldier goes, moved to his deepest
soul.  Impossible it seems to longer worship the Roman gods.  When
he has described to Aulus the Feast of Dedication, he repeats the
words he has heard in the Temple cloister, and says in deepest
seriousness:

"Most unearthly is the Man on whom I have looked to-day.  In his
speech a divine patience, kindness, and dignity combine.  As for
the words he spoke, I cannot tell their moving power.  The sayings
of our noblest Romans are feeble in the comparison.  Never have I
heard another speak as he has done about a future world.  Truly, an
unequaled Man is this new Teacher who is abroad in Judaea."

Sleep is of little consequence that night.  Is the word of the
augur at Brundisium beginning to be fulfilled?  In his tent Quintus
is wondering through the long hours if, among his people on the
Tiber, the Shepherd shall not find some sheep to whom he will give
the unending life.